\documentclass[11pt,a4paper]{article}
\usepackage[utf8]{inputenc}
\usepackage[T1]{fontenc}
\usepackage[catalan]{babel}
\usepackage{lmodern}
\usepackage[sfdefault,lining]{FiraSans}
\usepackage[a4paper,top=2.3cm,bottom=2.2cm,left=1.8cm,right=1.8cm,headheight=1.3cm,headsep=0.4cm,footskip=1cm]{geometry}
\usepackage{amsmath,amssymb}
\usepackage{icomma}
\usepackage{graphicx}
\usepackage[table]{xcolor}
\usepackage{tikz}
\usepackage{fontawesome5}
\usepackage{ccicons}
\usepackage{fancyhdr}
\usepackage{lastpage}
\usepackage{tasks}
\usepackage{needspace}
\usepackage[most]{tcolorbox}
\usepackage{array}
\usepackage[hidelinks,pdftitle={Solucions de la fitxa de la unitat 11: Cossos geomètrics. 3r d'ESO},pdfsubject={Matemàtiques 3r ESO},pdfkeywords={poliedres, fórmula d'Euler, poliedres regulars, prisma, piràmide, cilindre, con, esfera, àrea, volum, capacitat}]{hyperref}

\definecolor{navy}{HTML}{1B2748}
\definecolor{accent}{HTML}{E8603A}
\definecolor{teal}{HTML}{1F9E90}
\definecolor{soft}{HTML}{EEF1F8}
\definecolor{muted}{HTML}{5D6685}

\setlength{\parindent}{0pt}
\setlength{\parskip}{3pt}
\settasks{label=\alph*), label-format=\bfseries\color{navy}, label-width=1.3em, item-indent=1.8em, column-sep=0.6em, after-item-skip=5pt, before-skip=3pt, after-skip=4pt}

% ---------- capçalera i peu ----------
\newcommand{\llicencia}{\href{https://creativecommons.org/licenses/by-nc-sa/4.0/deed.ca}{\ccbyncsa}\ \ Llicència Creative Commons Reconeixement-NoComercial-CompartirIgual 4.0 Internacional (CC BY-NC-SA 4.0)}
\pagestyle{fancy}
\fancyhf{}
\fancyhead[L]{\raisebox{-0.35\height}{\includegraphics[height=0.95cm]{logo-gva.png}}}
\fancyhead[C]{\small\color{muted}Solucions de la fitxa · Unitat 11: Cossos geomètrics · 3r d'ESO}
\fancyhead[R]{\raisebox{-0.35\height}{\includegraphics[height=0.95cm]{logo-vermellar.png}}}
\fancyfoot[L]{\scriptsize\color{muted}\llicencia}
\fancyfoot[R]{\small\color{muted}\thepage}
\renewcommand{\headrulewidth}{0.4pt}
\renewcommand{\headrule}{\color{muted!40}\hrule height \headrulewidth}
\fancypagestyle{primera}{\fancyhf{}\renewcommand{\headrulewidth}{0pt}
  \fancyfoot[L]{\scriptsize\color{muted}\llicencia}
  \fancyfoot[R]{\small\color{muted}\thepage}}

% ---------- nivells i exercicis ----------
\newcommand{\estrelles}[1]{{\color{accent}\ifcase#1\or\faStar\faStar[regular]\faStar[regular]\or\faStar\faStar\faStar[regular]\or\faStar\faStar\faStar\fi}}
\newcounter{exer}
\newcommand{\ex}[1]{\par\needspace{4\baselineskip}\medskip\stepcounter{exer}\noindent{\color{navy}\bfseries\theexer.}\ \ {\small\estrelles{#1}}\ \ }

\newcommand{\seccio}[2]{\par\needspace{10\baselineskip}\bigskip
  \begin{tcolorbox}[colback=navy,colframe=navy,coltext=white,arc=2mm,boxrule=0pt,left=3mm,right=3mm,top=1.2mm,bottom=1.2mm]
  {\large\bfseries #1\quad #2}
  \end{tcolorbox}\vspace{-2pt}}

\newcolumntype{C}[1]{>{\centering\arraybackslash}p{#1}}
\usepackage{pifont}

% quadre de recordatori (només a la fitxa d'exercicis)
\newcommand{\recorda}[1]{\par\smallskip\begin{tcolorbox}[colback=accent!7,colframe=accent!55,arc=2mm,boxrule=0.5pt,left=2.5mm,right=2.5mm,top=1mm,bottom=1mm]\small{\color{accent}\faLightbulb[regular]}\ \textbf{Recorda:}\ #1\end{tcolorbox}}
\newcommand{\res}[1]{\boxed{#1}}
\newcommand{\rest}[1]{\fbox{#1}}

\def\fmneed{14}
% =====================================================================
% Capçalera FentMates de les fitxes: un color per a cada curs.
% Els colors són els de la pàgina inicial (LEVELS a assets/js/materials.js);
% el logo ve de _eines/logo/pdf/ (python3 _eines/logo/genera.py).
%
% Ús, al final del preàmbul de la unitat:
%   % =====================================================================
% Capçalera FentMates de les fitxes: un color per a cada curs.
% Els colors són els de la pàgina inicial (LEVELS a assets/js/materials.js);
% el logo ve de _eines/logo/pdf/ (python3 _eines/logo/genera.py).
%
% Ús, al final del preàmbul de la unitat:
%   % =====================================================================
% Capçalera FentMates de les fitxes: un color per a cada curs.
% Els colors són els de la pàgina inicial (LEVELS a assets/js/materials.js);
% el logo ve de _eines/logo/pdf/ (python3 _eines/logo/genera.py).
%
% Ús, al final del preàmbul de la unitat:
%   \input{../comu/fentmates.tex}
%   \fmcurs{1eso}
% i, al començament del document:
%   \capcalera{Matemàtiques}{1r ESO}{Fitxa d'exercicis · Unitat 2}{Divisibilitat}
% Opcional abans de \input: \def\fmneed{14} (espai mínim abans de cada secció).
% =====================================================================
\makeatletter
\def\fm@logodir{../../logo/pdf}
\providecommand{\fmneed}{10}
\def\fm@cursnom{}

% clau del curs, color i color del text damunt del color
\newcommand{\fm@curs}[3]{\@namedef{fm@color@#1}{#2}\@namedef{fm@text@#1}{#3}}
\fm@curs{1eso}{E6B84A}{navy}
\fm@curs{2eso}{3B82F6}{white}
\fm@curs{3eso}{FF7A50}{white}
\fm@curs{4eso-a}{A78BFA}{white}
\fm@curs{4eso-b}{6D28D9}{white}
\fm@curs{1batx-mi}{9B1C1C}{white}
\fm@curs{1batx-ccss}{EF6B6B}{white}
\fm@curs{2batx-mii}{1E6B47}{white}
\fm@curs{2batx-ccss}{2BB3A3}{white}

\newcommand{\fmcurs}[1]{%
  \@ifundefined{fm@color@#1}{\PackageError{fentmates}{Curs desconegut: #1}{}}{}%
  \edef\fm@tmp{\noexpand\definecolor{curs}{HTML}{\@nameuse{fm@color@#1}}%
               \noexpand\colorlet{curstext}{\@nameuse{fm@text@#1}}}\fm@tmp
  \colorlet{cursfosc}{curs!72!black}%
  \def\fm@logo{\fm@logodir/vertical-#1.pdf}%
  \def\fm@marca{\fm@logodir/marca-#1.pdf}}

% ---------- capçalera de la primera pàgina ----------
% #1 matèria (petit), #2 curs, #3 tipus de document i unitat, #4 títol
\newcommand{\capcalera}[4]{%
  \gdef\fm@cursnom{#2}%
  \par\noindent
  \begin{tcolorbox}[colback=curs,colframe=curs,boxrule=0pt,arc=2.2mm,width=3.5cm,height=2.15cm,
      valign=center,left=3mm,right=2mm,top=1mm,bottom=1mm,nobeforeafter,box align=center]
    \color{curstext}{\footnotesize #1}\par\vspace{1pt}{\fontsize{19}{21}\selectfont\bfseries #2\par}%
  \end{tcolorbox}\hspace{5mm}%
  \begin{minipage}[c]{\dimexpr\linewidth-3.5cm-5mm-2.9cm\relax}
    {\color{cursfosc}\small\bfseries\MakeUppercase{#3}}\par\vspace{2pt}
    {\color{navy}\fontsize{20}{23}\selectfont\bfseries #4\par}
  \end{minipage}\hfill
  \begin{minipage}[c]{2.6cm}\raggedleft\includegraphics[width=2.6cm]{\fm@logo}\end{minipage}\par
  \vspace{2pt}}

% ---------- barres de secció del color del curs ----------
\renewcommand{\seccio}[2]{\par\needspace{\fmneed\baselineskip}\bigskip
  \begin{tcolorbox}[colback=curs,colframe=curs,coltext=curstext,arc=2mm,boxrule=0pt,left=3mm,right=3mm,top=1.2mm,bottom=1.2mm]
  {\large\bfseries #1\quad #2}
  \end{tcolorbox}\vspace{-2pt}}

% ---------- capçalera i peu de les pàgines ----------
% primera pàgina: logos institucionals al peu
\fancypagestyle{primera}{\fancyhf{}\renewcommand{\headrulewidth}{0pt}
  \fancyfoot[L]{\raisebox{-0.32\height}{\includegraphics[height=0.78cm]{logo-gva.png}}\hspace{5mm}%
    \raisebox{-0.32\height}{\includegraphics[height=0.78cm]{logo-vermellar.png}}}
  \fancyfoot[R]{\scriptsize\color{muted}\href{https://creativecommons.org/licenses/by-nc-sa/4.0/deed.ca}{\ccbyncsa}\ \ CC BY-NC-SA 4.0 · fentmates.org\qquad\small\thepage}}
% la resta: el símbol a l'esquerra i el curs a la dreta, amb el títol al centre com abans
\fancyhead[L]{\raisebox{-0.3\height}{\includegraphics[height=0.62cm]{\fm@marca}}}
\fancyhead[R]{\small\bfseries\color{cursfosc}\fm@cursnom}
\renewcommand{\headrule}{\color{curs!45}\hrule height \headrulewidth}
\makeatother

%   \fmcurs{1eso}
% i, al començament del document:
%   \capcalera{Matemàtiques}{1r ESO}{Fitxa d'exercicis · Unitat 2}{Divisibilitat}
% Opcional abans de \input: \def\fmneed{14} (espai mínim abans de cada secció).
% =====================================================================
\makeatletter
\def\fm@logodir{../../logo/pdf}
\providecommand{\fmneed}{10}
\def\fm@cursnom{}

% clau del curs, color i color del text damunt del color
\newcommand{\fm@curs}[3]{\@namedef{fm@color@#1}{#2}\@namedef{fm@text@#1}{#3}}
\fm@curs{1eso}{E6B84A}{navy}
\fm@curs{2eso}{3B82F6}{white}
\fm@curs{3eso}{FF7A50}{white}
\fm@curs{4eso-a}{A78BFA}{white}
\fm@curs{4eso-b}{6D28D9}{white}
\fm@curs{1batx-mi}{9B1C1C}{white}
\fm@curs{1batx-ccss}{EF6B6B}{white}
\fm@curs{2batx-mii}{1E6B47}{white}
\fm@curs{2batx-ccss}{2BB3A3}{white}

\newcommand{\fmcurs}[1]{%
  \@ifundefined{fm@color@#1}{\PackageError{fentmates}{Curs desconegut: #1}{}}{}%
  \edef\fm@tmp{\noexpand\definecolor{curs}{HTML}{\@nameuse{fm@color@#1}}%
               \noexpand\colorlet{curstext}{\@nameuse{fm@text@#1}}}\fm@tmp
  \colorlet{cursfosc}{curs!72!black}%
  \def\fm@logo{\fm@logodir/vertical-#1.pdf}%
  \def\fm@marca{\fm@logodir/marca-#1.pdf}}

% ---------- capçalera de la primera pàgina ----------
% #1 matèria (petit), #2 curs, #3 tipus de document i unitat, #4 títol
\newcommand{\capcalera}[4]{%
  \gdef\fm@cursnom{#2}%
  \par\noindent
  \begin{tcolorbox}[colback=curs,colframe=curs,boxrule=0pt,arc=2.2mm,width=3.5cm,height=2.15cm,
      valign=center,left=3mm,right=2mm,top=1mm,bottom=1mm,nobeforeafter,box align=center]
    \color{curstext}{\footnotesize #1}\par\vspace{1pt}{\fontsize{19}{21}\selectfont\bfseries #2\par}%
  \end{tcolorbox}\hspace{5mm}%
  \begin{minipage}[c]{\dimexpr\linewidth-3.5cm-5mm-2.9cm\relax}
    {\color{cursfosc}\small\bfseries\MakeUppercase{#3}}\par\vspace{2pt}
    {\color{navy}\fontsize{20}{23}\selectfont\bfseries #4\par}
  \end{minipage}\hfill
  \begin{minipage}[c]{2.6cm}\raggedleft\includegraphics[width=2.6cm]{\fm@logo}\end{minipage}\par
  \vspace{2pt}}

% ---------- barres de secció del color del curs ----------
\renewcommand{\seccio}[2]{\par\needspace{\fmneed\baselineskip}\bigskip
  \begin{tcolorbox}[colback=curs,colframe=curs,coltext=curstext,arc=2mm,boxrule=0pt,left=3mm,right=3mm,top=1.2mm,bottom=1.2mm]
  {\large\bfseries #1\quad #2}
  \end{tcolorbox}\vspace{-2pt}}

% ---------- capçalera i peu de les pàgines ----------
% primera pàgina: logos institucionals al peu
\fancypagestyle{primera}{\fancyhf{}\renewcommand{\headrulewidth}{0pt}
  \fancyfoot[L]{\raisebox{-0.32\height}{\includegraphics[height=0.78cm]{logo-gva.png}}\hspace{5mm}%
    \raisebox{-0.32\height}{\includegraphics[height=0.78cm]{logo-vermellar.png}}}
  \fancyfoot[R]{\scriptsize\color{muted}\href{https://creativecommons.org/licenses/by-nc-sa/4.0/deed.ca}{\ccbyncsa}\ \ CC BY-NC-SA 4.0 · fentmates.org\qquad\small\thepage}}
% la resta: el símbol a l'esquerra i el curs a la dreta, amb el títol al centre com abans
\fancyhead[L]{\raisebox{-0.3\height}{\includegraphics[height=0.62cm]{\fm@marca}}}
\fancyhead[R]{\small\bfseries\color{cursfosc}\fm@cursnom}
\renewcommand{\headrule}{\color{curs!45}\hrule height \headrulewidth}
\makeatother

%   \fmcurs{1eso}
% i, al començament del document:
%   \capcalera{Matemàtiques}{1r ESO}{Fitxa d'exercicis · Unitat 2}{Divisibilitat}
% Opcional abans de \input: \def\fmneed{14} (espai mínim abans de cada secció).
% =====================================================================
\makeatletter
\def\fm@logodir{../../logo/pdf}
\providecommand{\fmneed}{10}
\def\fm@cursnom{}

% clau del curs, color i color del text damunt del color
\newcommand{\fm@curs}[3]{\@namedef{fm@color@#1}{#2}\@namedef{fm@text@#1}{#3}}
\fm@curs{1eso}{E6B84A}{navy}
\fm@curs{2eso}{3B82F6}{white}
\fm@curs{3eso}{FF7A50}{white}
\fm@curs{4eso-a}{A78BFA}{white}
\fm@curs{4eso-b}{6D28D9}{white}
\fm@curs{1batx-mi}{9B1C1C}{white}
\fm@curs{1batx-ccss}{EF6B6B}{white}
\fm@curs{2batx-mii}{1E6B47}{white}
\fm@curs{2batx-ccss}{2BB3A3}{white}

\newcommand{\fmcurs}[1]{%
  \@ifundefined{fm@color@#1}{\PackageError{fentmates}{Curs desconegut: #1}{}}{}%
  \edef\fm@tmp{\noexpand\definecolor{curs}{HTML}{\@nameuse{fm@color@#1}}%
               \noexpand\colorlet{curstext}{\@nameuse{fm@text@#1}}}\fm@tmp
  \colorlet{cursfosc}{curs!72!black}%
  \def\fm@logo{\fm@logodir/vertical-#1.pdf}%
  \def\fm@marca{\fm@logodir/marca-#1.pdf}}

% ---------- capçalera de la primera pàgina ----------
% #1 matèria (petit), #2 curs, #3 tipus de document i unitat, #4 títol
\newcommand{\capcalera}[4]{%
  \gdef\fm@cursnom{#2}%
  \par\noindent
  \begin{tcolorbox}[colback=curs,colframe=curs,boxrule=0pt,arc=2.2mm,width=3.5cm,height=2.15cm,
      valign=center,left=3mm,right=2mm,top=1mm,bottom=1mm,nobeforeafter,box align=center]
    \color{curstext}{\footnotesize #1}\par\vspace{1pt}{\fontsize{19}{21}\selectfont\bfseries #2\par}%
  \end{tcolorbox}\hspace{5mm}%
  \begin{minipage}[c]{\dimexpr\linewidth-3.5cm-5mm-2.9cm\relax}
    {\color{cursfosc}\small\bfseries\MakeUppercase{#3}}\par\vspace{2pt}
    {\color{navy}\fontsize{20}{23}\selectfont\bfseries #4\par}
  \end{minipage}\hfill
  \begin{minipage}[c]{2.6cm}\raggedleft\includegraphics[width=2.6cm]{\fm@logo}\end{minipage}\par
  \vspace{2pt}}

% ---------- barres de secció del color del curs ----------
\renewcommand{\seccio}[2]{\par\needspace{\fmneed\baselineskip}\bigskip
  \begin{tcolorbox}[colback=curs,colframe=curs,coltext=curstext,arc=2mm,boxrule=0pt,left=3mm,right=3mm,top=1.2mm,bottom=1.2mm]
  {\large\bfseries #1\quad #2}
  \end{tcolorbox}\vspace{-2pt}}

% ---------- capçalera i peu de les pàgines ----------
% primera pàgina: logos institucionals al peu
\fancypagestyle{primera}{\fancyhf{}\renewcommand{\headrulewidth}{0pt}
  \fancyfoot[L]{\raisebox{-0.32\height}{\includegraphics[height=0.78cm]{logo-gva.png}}\hspace{5mm}%
    \raisebox{-0.32\height}{\includegraphics[height=0.78cm]{logo-vermellar.png}}}
  \fancyfoot[R]{\scriptsize\color{muted}\href{https://creativecommons.org/licenses/by-nc-sa/4.0/deed.ca}{\ccbyncsa}\ \ CC BY-NC-SA 4.0 · fentmates.org\qquad\small\thepage}}
% la resta: el símbol a l'esquerra i el curs a la dreta, amb el títol al centre com abans
\fancyhead[L]{\raisebox{-0.3\height}{\includegraphics[height=0.62cm]{\fm@marca}}}
\fancyhead[R]{\small\bfseries\color{cursfosc}\fm@cursnom}
\renewcommand{\headrule}{\color{curs!45}\hrule height \headrulewidth}
\makeatother

\fmcurs{3eso}

\begin{document}
\thispagestyle{primera}

\capcalera{Matemàtiques}{3r ESO}{Solucions de la fitxa · Unitat 11}{Cossos geomètrics}
\vspace{8pt}

\begin{tcolorbox}[colback=soft,colframe=soft,arc=2mm,left=3mm,right=3mm,top=1.5mm,bottom=1.5mm]
\small Intenta resoldre cada exercici abans de mirar-ne la solució. Si no t'ix el mateix resultat, repassa els passos o consulta els apunts interactius de la unitat: hi trobaràs cossos que pots girar i eines que calculen àrees i volums pas a pas. Els resultats aproximats estan arrodonits a les centèsimes amb la tecla $\pi$ de la calculadora.
\end{tcolorbox}

% =====================================================
\seccio{1}{Poliedres. Fórmula d'Euler. Poliedres regulars}

\ex{1} \begin{center}\renewcommand{\arraystretch}{1.25}\begin{tabular}{|l|C{1.6cm}|C{1.6cm}|C{1.6cm}|l|}\hline Poliedre & $C$ & $V$ & $A$ & $C+V=A+2$ \\\hline prisma triangular & 5 & 6 & 9 & $11=11$ \\\hline piràmide pentagonal & 6 & 6 & 10 & $12=12$ \\\hline prisma hexagonal & 8 & 12 & 18 & $20=20$ \\\hline piràmide octogonal & 9 & 9 & 16 & $18=18$ \\\hline cub & 6 & 8 & 12 & $14=14$ \\\hline\end{tabular}\end{center}

\ex{1} \begin{tasks}(3)
  \task $8+12=A+2\Rightarrow\res{A=18}$
  \task $C+10=17\Rightarrow\res{C=7}$
  \task $12+V=32\Rightarrow\res{V=20}$
\end{tasks}

\ex{1} \begin{tasks}(3)
  \task \rest{sí}
  \task no (rectangles diferents)
  \task \rest{sí}
  \task no (cares de dos tipus)
  \task \rest{sí}
  \task \rest{sí}
\end{tasks}Un poliedre regular té totes les cares polígons regulars iguals i el mateix nombre de cares a cada vèrtex.

\ex{2} $3n=24\Rightarrow n=8$: base \rest{octogonal}; \ $\res{C=10}$, \ $\res{V=16}$ \ ($10+16=24+2$).

\ex{2} $n+1=11\Rightarrow n=10$: base \rest{decagonal}; \ $\res{C=11}$, \ $\res{A=20}$ \ ($11+11=20+2$).

\ex{2} a) $7+7=14=12+2$: \rest{sí} (per exemple, una piràmide hexagonal). \ b) $5+6=11\neq12$: \rest{no}.

\ex{3} Quadrangulars: $\res{C=8,\ V=6,\ A=12}$ (és un octaedre si les cares són equilàters). \ Pentagonals: $\res{C=10,\ V=7,\ A=15}$. \ En tots dos casos $C+V=A+2$.

\ex{3} Cada angle d'un hexàgon regular fa $120^\circ$. En un vèrtex s'hi han d'ajuntar com a mínim 3 cares, i $3\cdot120^\circ=360^\circ$: les cares quedarien \rest{planes} i no tancarien cap cos.


% =====================================================
\seccio{2}{Prismes: elements i àrea}

\ex{1} \begin{tasks}(2)
  \task $6\cdot4^2=\res{96\text{ cm}^2}$
  \task $2(6\cdot3+6\cdot4+3\cdot4)=2\cdot54=\res{108\text{ cm}^2}$
\end{tasks}

\ex{1} $D=\sqrt{144+16+9}=\sqrt{169}=\res{13\text{ cm}}$.

\ex{1} $A_L=20\cdot12=\res{240\text{ cm}^2}$; \ $A_T=240+2\cdot25=\res{290\text{ cm}^2}$.

\ex{2} $A_L=12\cdot10=120$; altura de la base $\sqrt{16-4}=\sqrt{12}\approx3{,}46$; $A_{\text{base}}=\frac{4\cdot\sqrt{12}}{2}\approx6{,}93$. \ $A_T\approx120+2\cdot6{,}93=\res{133{,}86\text{ cm}^2}$.

\ex{2} $A_L=36\cdot8=288$; \ $a_b=\sqrt{36-9}=\sqrt{27}\approx5{,}20$; \ $A_{\text{base}}=\frac{36\cdot\sqrt{27}}{2}\approx93{,}53$. \ $A_T\approx288+187{,}06=\res{475{,}06\text{ cm}^2}$.

\ex{2} Base $40\cdot30=1200$; \ cares $2\cdot40\cdot20+2\cdot30\cdot20=2800$. \ Total: $\res{4000\text{ cm}^2}$.

\ex{3} $D^2=3l^2\Rightarrow3l^2=48\Rightarrow\res{l=4\text{ cm}}$; \ $A_T=6\cdot16=\res{96\text{ cm}^2}$.

\ex{3} Hipotenusa: $\sqrt{36+64}=10$. \ $A_L=(6+8+10)\cdot15=360$; \ $A_{\text{base}}=\frac{6\cdot8}{2}=24$. \ $A_T=360+48=\res{408\text{ cm}^2}$.


% =====================================================
\seccio{3}{Piràmides i troncs de piràmide}

\ex{1} $a_b=5$; \ $a_p=\sqrt{144+25}=\res{13\text{ cm}}$. \ $A_T=\frac{40\cdot13}{2}+100=\res{360\text{ cm}^2}$.

\ex{1} $A_L=\frac{24\cdot9}{2}=\res{108\text{ cm}^2}$.

\ex{2} $h=\sqrt{5^2-3^2}=\res{4\text{ cm}}$; \ $A_T=\frac{24\cdot5}{2}+36=\res{96\text{ cm}^2}$.

\ex{2} $a_p=\sqrt{36+16}=\sqrt{52}\approx\res{7{,}21\text{ cm}}$; \ aresta: $\sqrt{a_p^2+4^2}=\sqrt{68}\approx\res{8{,}25\text{ cm}}$; \ $A_T=\frac{32\cdot\sqrt{52}}{2}+64\approx\res{179{,}38\text{ cm}^2}$.

\ex{2} Cada cara és un triangle equilàter d'altura $\sqrt{36-9}=\sqrt{27}\approx5{,}20$: \ $A=4\cdot\frac{6\cdot\sqrt{27}}{2}\approx\res{62{,}35\text{ cm}^2}$.

\ex{3} $a_p=\sqrt{4^2+3^2}=\res{5\text{ cm}}$ \ (catets: l'altura i $\frac{12-6}{2}=3$). \ $A_T=\frac{(48+24)\cdot5}{2}+144+36=\res{360\text{ cm}^2}$.

\ex{3} $a_p=\sqrt{146^2+115^2}=\sqrt{34\,541}\approx\res{185{,}85\text{ m}}$; \ $A_L=\frac{920\cdot185{,}85}{2}\approx\res{85\,492\text{ m}^2}$.


% =====================================================
\seccio{4}{Cilindre i con}

\ex{1} $A_L=2\pi\cdot3\cdot10=60\pi\approx\res{188{,}50\text{ cm}^2}$; \ $A_T=60\pi+18\pi=78\pi\approx\res{245{,}04\text{ cm}^2}$.

\ex{1} $g=\sqrt{144+25}=\res{13\text{ cm}}$; \ $A_L=\pi\cdot5\cdot13=65\pi\approx\res{204{,}20\text{ cm}^2}$; \ $A_T=65\pi+25\pi=90\pi\approx\res{282{,}74\text{ cm}^2}$.

\ex{1} $h=\sqrt{100-36}=\res{8\text{ cm}}$; \ $A_T=60\pi+36\pi=96\pi\approx\res{301{,}59\text{ cm}^2}$.

\ex{2} $r=4$: \ $A_T=2\pi\cdot4\cdot12+2\pi\cdot16=96\pi+32\pi=128\pi\approx\res{402{,}12\text{ cm}^2}$.

\ex{2} $g=\sqrt{81+144}=15$; \ només la superfície lateral: $A_L=\pi\cdot9\cdot15=135\pi\approx\res{424{,}12\text{ cm}^2}$.

\ex{2} $2\pi r\cdot10=120\pi\Rightarrow\res{r=6\text{ cm}}$; \ $A_T=120\pi+72\pi=192\pi\approx\res{603{,}19\text{ cm}^2}$.

\ex{3} L'arc del sector fa $2\pi\cdot4=8\pi$ i la circumferència sencera de radi 12 faria $24\pi$: és $\frac{8\pi}{24\pi}=\frac13$ del cercle. \ Angle: $\frac{360^\circ}{3}=\res{120^\circ}$.

\ex{3} $A=12\cdot2\pi\cdot0{,}4\cdot5=48\pi\approx150{,}80\text{ m}^2$; \ $\frac{150{,}80}{4}\approx\res{37{,}70\text{ L}}$.


% =====================================================
\seccio{5}{L'esfera. La Terra}

\ex{1} \begin{tasks}(3)
  \task $64\pi\approx\res{201{,}06\text{ cm}^2}$
  \task $100\pi\approx\res{314{,}16\text{ cm}^2}$
  \task $\pi\approx\res{3{,}14\text{ m}^2}$
\end{tasks}

\ex{2} $4\pi r^2=144\pi\Rightarrow r^2=36\Rightarrow\res{r=6\text{ cm}}$.

\ex{2} Mitja esfera: $\frac{4\pi\cdot25}{2}=50\pi\approx\res{157{,}08\text{ m}^2}$.

\ex{2} $40^\circ+10^\circ=\res{50^\circ}$; \ $\frac{40\,000\cdot50}{360}\approx\res{5555{,}56\text{ km}}$.

\ex{2} Cada $15^\circ$ és una hora: \ $45^\circ$ E $\to$ 3 hores més: \res{15 h}; \ $30^\circ$ O $\to$ 2 hores menys: \res{10 h}.

\ex{3} $\frac{4\pi\cdot6371^2}{4\pi\cdot1737^2}=\left(\frac{6371}{1737}\right)^2\approx\res{13{,}45}$ vegades.


% =====================================================
\seccio{6}{Volum. Prismes i cilindres}

\ex{1} \begin{tasks}(3)
  \task $\res{3200\text{ L}}$
  \task $\res{0{,}45\text{ L}}$
  \task $\res{2500\text{ cm}^3}$
  \task $\res{750\text{ mL}}$
  \task $\res{12\text{ m}^3}$
  \task $\res{0{,}5\text{ m}^3}$
\end{tasks}

\ex{1} \begin{tasks}(3)
  \task $6^3=\res{216\text{ cm}^3}$
  \task $\res{120\text{ cm}^3}$
  \task $160\pi\approx\res{502{,}65\text{ cm}^3}$
\end{tasks}

\ex{1} $V=\frac{5\cdot12}{2}\cdot20=\res{600\text{ cm}^3}$.

\ex{2} $a_b=\sqrt{4-1}=\sqrt3$; \ $A_{\text{base}}=\frac{12\cdot\sqrt3}{2}=6\sqrt3\approx10{,}39$; \ $V\approx\res{51{,}96\text{ cm}^3}$.

\ex{2} $V=72\,000\text{ cm}^3=\res{72\text{ L}}$; \ $\frac34\cdot72=\res{54\text{ L}}$.

\ex{2} $V=\pi\cdot1{,}5^2\cdot4=9\pi\approx28{,}27\text{ m}^3\approx\res{28\,274\text{ L}}$.

\ex{2} Base (el lateral de la piscina): $\frac{(1+3)\cdot25}{2}=50\text{ m}^2$; \ $V=50\cdot10=\res{500\text{ m}^3}$.

\ex{3} $\pi\cdot16\cdot h=500\Rightarrow h=\frac{500}{16\pi}\approx\res{9{,}95\text{ cm}}$.

\ex{3} $V=24\cdot11{,}5\cdot5-3\cdot\pi\cdot1{,}5^2\cdot5=1380-33{,}75\pi\approx\res{1273{,}97\text{ cm}^3}$.


% =====================================================
\seccio{7}{Volum de piràmides, cons i esferes}

\ex{1} \begin{tasks}(3)
  \task $\frac{36\cdot10}{3}=\res{120\text{ cm}^3}$
  \task $21\pi\approx\res{65{,}97\text{ cm}^3}$
  \task $288\pi\approx\res{904{,}78\text{ cm}^3}$
\end{tasks}

\ex{1} $r=10$: \ $V=\frac43\pi\cdot1000\approx\res{4188{,}79\text{ cm}^3}\approx\res{4{,}19\text{ L}}$.

\ex{2} $h=\sqrt{13^2-5^2}=\res{12\text{ cm}}$; \ $V=\frac{100\cdot12}{3}=\res{400\text{ cm}^3}$.

\ex{2} $h=\sqrt{100-36}=8$; \ $V=\frac{\pi\cdot36\cdot8}{3}=96\pi\approx\res{301{,}59\text{ cm}^3}$.

\ex{2} $V=\frac43\pi\cdot11^3\approx5575{,}28\text{ cm}^3\approx\res{5{,}58\text{ L}}$.

\ex{2} $a_b=\sqrt{12}$; \ $A_{\text{base}}=\frac{24\cdot\sqrt{12}}{2}\approx41{,}57$; \ $V=\frac{A_{\text{base}}\cdot9}{3}\approx\res{124{,}71\text{ cm}^3}$.

\ex{3} $V=\frac12\cdot\frac43\pi\cdot12^3=1152\pi\approx3619{,}11\text{ cm}^3\approx\res{3{,}62\text{ L}}$.

\ex{3} $\frac13\pi\cdot36\cdot h=\frac43\pi\cdot216\Rightarrow12h=288\Rightarrow\res{h=24\text{ cm}}$ \ (quatre vegades el radi).


% =====================================================
\seccio{8}{Cossos compostos i problemes}

\ex{1} $V=6^3+\frac{36\cdot4}{3}=216+48=\res{264\text{ cm}^3}$.

\ex{2} $V=\pi\cdot9\cdot10+\frac{\pi\cdot9\cdot4}{3}=90\pi+12\pi=102\pi\approx\res{320{,}44\text{ m}^3}$.

\ex{2} $V=\pi\cdot4\cdot6+\frac43\pi\cdot8=24\pi+\frac{32\pi}{3}=\frac{104\pi}{3}\approx\res{108{,}91\text{ cm}^3}$; \ $A=2\pi\cdot2\cdot6+4\pi\cdot4=40\pi\approx\res{125{,}66\text{ cm}^2}$.

\ex{2} $V=\frac{\pi\cdot2{,}5^2\cdot12}{3}+\frac12\cdot\frac43\pi\cdot2{,}5^3=25\pi+\frac{125\pi}{12}=\frac{425\pi}{12}\approx\res{111{,}26\text{ cm}^3}$.

\ex{2} $V=\pi(6^2-5^2)\cdot50=550\pi\approx\res{1727{,}88\text{ cm}^3}$.

\ex{2} $V=1000-\pi\cdot9\cdot10=1000-90\pi\approx\res{717{,}26\text{ cm}^3}$.

\ex{3} Si el con menut fa $x$ d'altura: $\frac{x}{12}=\frac{x+24}{18}\Rightarrow x=48$. \ $V=\frac{\pi\cdot18^2\cdot72}{3}-\frac{\pi\cdot12^2\cdot48}{3}=7776\pi-2304\pi=5472\pi\approx17\,190{,}80\text{ cm}^3\approx\res{17{,}19\text{ L}}$.

\ex{3} El volum de la bola és $\frac43\pi\cdot27=36\pi$. \ $\pi\cdot36\cdot h=36\pi\Rightarrow\res{h=1\text{ cm}}$.

\ex{3} $V=\pi\cdot4\cdot5+\frac12\cdot\frac43\pi\cdot8=20\pi+\frac{16\pi}{3}=\frac{76\pi}{3}\approx\res{79{,}59\text{ m}^3}$. \ $\frac{79\,587}{50}\approx1591{,}7$ min $\approx\res{26{,}5\text{ hores}}$.


\end{document}
